\chapter{Introduction}
\section{Motivation and Clinical Relevance}

Magnetic resonance imaging (MRI) is an important diagnostic imaging technique that provides detailed anatomical information without exposing patients to ionizing radiation. Its ability to distinguish soft tissues makes it particularly valuable for investigating fetal brain development. In addition to conventional anatomical imaging, diffusion magnetic resonance imaging (dMRI) provides information about tissue microstructure by characterizing the diffusion of water molecules within biological tissues. This enables the investigation of developing neural structures and white matter organization, contributing to the understanding of fetal brain maturation and developmental abnormalities \autocite{malamateniou2013motion,snoussi2025haitch}.

Despite these advantages, obtaining reliable fetal diffusion MRI data remains challenging. Spontaneous fetal movement, maternal respiration, and magnetic field inhomogeneities can introduce artifacts that degrade image quality and affect subsequent analysis. Diffusion-weighted acquisitions are particularly sensitive to motion because movement during diffusion encoding can alter the measured signal. Furthermore, the relatively low signal-to-noise ratio of fetal diffusion images increases the difficulty of obtaining consistent measurements. Such disturbances may lead to signal loss, geometric distortions, and inconsistencies between acquired images, potentially compromising the accuracy of diffusion-related measurements and the interpretation of fetal brain microstructure \autocite{malamateniou2013motion,snoussi2025haitch}.

Identifying artifact-corrupted acquisitions is therefore an important aspect of fetal dMRI quality control. Conventional retrospective assessment allows corrupted images to be identified after acquisition, while automated quality-control approaches offer opportunities for evaluating data with less manual intervention. The High-Efficiency Real-Time Motion Quantification and Re-Acquisition (HERON) framework was developed to support automated fetal diffusion MRI acquisition through motion assessment, image-quality evaluation, and reacquisition of affected diffusion volumes. Its development demonstrates the potential of incorporating quality assessment directly into fetal MRI acquisition workflows \autocite{avilesverdera2025heron}.

Existing methods for automated fetal MRI quality assessment employ different strategies, including image-derived quality measures, motion-related criteria, and machine learning. For example, FetMRQC combines image-quality metrics with a Random Forest model to predict expert quality ratings in fetal brain MRI. Although developed for anatomical T2-weighted images, this approach demonstrates how learned relationships between image characteristics and quality annotations can support automated assessment \autocite{sanchez2024fetmrqc}. These developments motivate further investigation of learning-based methods that can recognize artifact-related image characteristics directly from fetal diffusion MRI data.

Within this context, the present thesis investigates two deep learning-based approaches for automated artifact detection in low-field fetal dMRI. The first approach uses three-dimensional (3D) image information to assess complete diffusion dynamics, while the second examines individual two-dimensional (2D) slices to investigate artifact detection at a finer processing level. Both approaches combine automatic fetal-brain localization with image-quality classification, enabling the evaluation of artifact-detection performance using different spatial representations.

The research is further motivated by the possibility of incorporating automated artifact assessment into acquisition-oriented processing workflows. As direct access to an MRI scanner was unavailable for experimental integration, the developed methods were investigated through retrospective evaluation and offline scanner-simulation experiments. This provides a basis for studying the capabilities, limitations, and computational characteristics of the proposed approaches without assuming that real-time clinical performance has been established.

\section{Problem Statement}

The central challenge is to distinguish usable fetal diffusion MRI measurements from artifact-corrupted measurements when corruption may vary between individual slices and complete volumes \autocite{malamateniou2013motion,snoussi2025haitch}. This creates uncertainty about which image representation provides sufficient evidence for a reliable quality decision.

Volumetric analysis preserves relationships between slices but requires the dynamic to be available before assessment. Slice-wise analysis can examine individual measurements earlier, but may miss degradation that becomes evident only across a volume. A direct comparison is needed rather than assuming either representation is preferable.

Both approaches depend on reliable fetal-brain localization before classification. Segmentation errors can exclude relevant anatomy or prevent crop generation, so classification results must be interpreted in relation to the quality of the preceding localization stage.

The available artifact annotations are primarily assigned to complete diffusion dynamics. Using those labels as targets for individual slices creates weak supervision: not every slice from an artifact-positive dynamic necessarily contains visible corruption. This mismatch must be acknowledged in model evaluation.

Slice-wise predictions also require a rule for deciding whether the complete dynamic is acceptable, including how to treat missing predictions caused by segmentation failure.

Scanner integration was not available in this study. Consequently, processing-time measurements and offline scanner-simulation experiments are used to investigate computational behavior and sequential decisions; they cannot establish clinical real-time performance.

The thesis therefore evaluates how image representation, brain localization, annotation granularity, and dynamic-level decision rules influence learned artifact detection, within the limits of retrospective data and offline simulations.

\section{Research Objectives and Research Questions}

The objective is to establish the capabilities and limitations of two deep learning approaches to artifact detection in low-field fetal diffusion MRI, and to determine what conclusions can be supported by the available data and offline evaluations.

The experiments examine 3D dynamic-level assessment and 2D slice-wise assessment as complementary approaches, without assuming that either representation is inherently superior.

The study also examines segmentation-dependent classification, uncertainty from volume-level labels used on slices, dynamic-level aggregation, and computational behavior under simulated acquisition workflows.

The research is guided by the following questions:

\noindent\textbf{RQ1: How effectively can a 3D deep learning pipeline identify artifact-corrupted fetal diffusion MRI dynamics using volumetric image information?}

This question addresses the development and evaluation of fetal-brain segmentation, volumetric region-of-interest extraction, and dynamic-level artifact classification. It investigates the performance and limitations of the nnU-Net v2 and 3D ResNet framework.

\noindent\textbf{RQ2: How effectively can a 2D deep learning pipeline perform fetal-brain localization and artifact classification using individual diffusion MRI slices?}

This question examines the performance of MobileNetV3-U-Net segmentation and EfficientNet-B0 classification, including the influence of reference-mask and predicted-mask cropping, segmentation failures, and uncertainty in slice-level quality annotations.

\noindent\textbf{RQ3: What differences can be observed between volumetric and slice-wise artifact detection in terms of classification performance and the information available for quality assessment?}

This question evaluates the recorded differences between the 3D and 2D approaches, particularly their use of spatial context, classification errors, and processing requirements. The comparison considers differences in experimental design and the available evaluation evidence.

\noindent\textbf{RQ4: How do slice-level prediction aggregation and segmentation-failure handling influence artifact detection at the dynamic level?}

This question investigates different strategies for combining slice-level classification outputs into quality decisions for complete diffusion dynamics. It examines the influence of aggregation methods, decision thresholds, and unsuccessful brain-region extraction on classification performance.

\noindent\textbf{RQ5: What can offline scanner-simulation experiments and computational measurements establish about the potential use of the developed approaches in acquisition-oriented quality control?}

This question examines the simulated processing workflows, computational latency, and sequential quality-assessment experiments. The evaluation focuses on the evidence provided by offline experiments rather than direct real-time scanner integration.

The original research proposal additionally identified evaluation across acquisition parameters, gestational age, and fetal pathology, together with potential generalization across different MRI field strengths. These aspects form part of the initial research context; however, the experimental scope is determined by the dataset ultimately provided and the analyses documented in the completed work. Cross-field generalization and prospective scanner integration are not treated as demonstrated outcomes of this thesis.

These questions define the evaluation scope used throughout the thesis.

\section{Main Contributions}

The work contributes implemented pipelines and experimental analyses addressing the research questions above. The contributions are distinguished by the processing systems and evaluations performed, rather than by claims of prospective clinical deployment.

The main contributions of this work are summarized as follows:

\begin{enumerate}
    \item \textbf{Development and evaluation of a volumetric artifact-detection pipeline.}

    A three-dimensional processing framework was developed by combining nnU-Net v2 for automatic fetal-brain segmentation with a 3D ResNet classifier for identifying artifact-corrupted diffusion dynamics. The framework incorporates volumetric brain-region extraction and dynamic-level quality classification. Its development was preceded by preliminary experiments involving custom 3D U-Net segmentation, handcrafted image features, and Random Forest classification. The resulting pipeline was evaluated to investigate the capabilities and limitations of artifact detection using complete diffusion volumes.

    \item \textbf{Development and evaluation of a slice-wise artifact-detection pipeline.}

    A complementary two-dimensional framework was developed using MobileNetV3-U-Net for fetal-brain segmentation and EfficientNet-B0 for artifact classification. The approach processes individual diffusion MRI slices and extracts brain-centered regions before classification. Its evaluation includes segmentation performance, classification reliability, and the influence of reference-mask and predicted-mask cropping. This provides an experimental basis for examining artifact detection when classification is performed using individual slices rather than complete volumes.

    \item \textbf{Investigation of segmentation reliability and its influence on artifact classification.}

    Both pipelines incorporate automatic fetal-brain localization as an essential processing stage. The research evaluates the relationship between segmentation quality, brain-region extraction, and subsequent artifact classification. In particular, the 2D experiments examine unsuccessful crop generation, differences between reference and predicted brain regions, and the effects of segmentation failures on classification decisions. The investigation also considers the limitations introduced when dynamic-level artifact annotations are used as slice-level classification targets.

    \item \textbf{Evaluation of dynamic-level decision strategies and comparison of the two processing approaches.}

    The research investigates methods for combining individual-slice artifact predictions into a quality decision for an entire diffusion dynamic. Different prediction aggregation strategies, classification thresholds, and segmentation-failure-aware decision rules are examined. A matched experimental comparison of the 3D and 2D pipelines is also conducted to investigate differences in artifact-detection behavior and classification errors, while accounting for limitations in the available evaluation evidence.

    \item \textbf{Investigation of computational performance through offline scanner-simulation experiments.}

    The developed approaches were further investigated through computational measurements and offline acquisition-workflow simulations. These experiments examine the processing characteristics of volumetric and slice-wise artifact assessment, including sequential evaluation and the potential for earlier quality decisions. Because direct MRI scanner access was unavailable, the investigation is limited to retrospective and simulated workflows and does not establish real-time clinical integration.
\end{enumerate}

The experimental evidence and its limitations are reported in Chapters 5 and 6.

\section{Organization of the Thesis}

This thesis is organized into seven chapters. Chapter 2 reviews the existing literature on fetal diffusion MRI, image artifacts, and automated quality assessment. Chapter 3 introduces the theoretical foundations of diffusion MRI, deep learning, image segmentation, and classification. Chapter 4 describes the dataset, preprocessing procedures, and methodology used to develop the 3D and 2D artifact-detection pipelines. Chapter 5 presents the experimental setup, evaluation procedures, and results obtained from both approaches, including their comparative analysis. Chapter 6 discusses the findings, methodological limitations, and implications of the results. Finally, Chapter 7 summarizes the main conclusions and outlines directions for future research.
